\section{The Simulator and Its Fidelity}
\label{sec:simulator}

\paraq{How can we choose a hybrid configuration without running full training on every one?}
Sweeping every $(k, \text{stream}, \text{graph}, \text{depth}, \text{hidden})$
configuration on multiple GPUs is prohibitively expensive. Instead we build a
\textbf{cost-model simulator} that predicts the expected per-epoch performance of a
configuration using \emph{few resources}---on the order of a \textbf{single GPU}---and
\emph{without finishing a full training run}. The simulator sweeps the split-point $k$
($k{=}0$ full split \ldots $k{=}L$ full data parallelism) and, for each, predicts the
\textbf{compute}, \textbf{communication}, and \textbf{data-loading} cost
(\texttt{get\_best\_strategy.py} $\rightarrow$ \texttt{simulation/hybrid\_dp\_sim.py}).

\paraq{What is the cheap-prediction philosophy here?}
The simulator times the model's per-layer compute once on one GPU, and predicts
communication and data-loading from feature \emph{volume} $\div$ hardware bandwidth,
rather than measuring an end-to-end multi-GPU run per configuration. A measured phase at
one point is projected across configurations analytically, so the whole configuration
space is covered from a handful of cheap measurements. We state the approximation
honestly: each layer's compute is \emph{measured} (a real forward$+$backward on dummy
inputs), not fitted to a closed-form linear cost, so the estimate inherits the true
GEMM/aggregation scaling; and the single-GPU timing holds \emph{no optimizer or
gradient-accumulation state} across layers---it charges the per-layer step cost and omits
the one-off optimizer update and parameter all-reduce, which are negligible beside the
per-layer shuffle in the regimes we study.

\subsection{Validation Against the DGL Baseline}
\label{sec:sim-baseline}

\paraq{Before trusting the simulator on hybrids, what must it first reproduce?}
A note on validation order: the simulator must first accurately reproduce the
\textbf{DGL baseline} before we rely on any hybrid prediction. DGL is the simplest
case---replicated graph, no partition-aware shuffle---so if the compute and data-loading
models cannot match DGL, no hybrid prediction is trustworthy. \placeholder{} (DGL
predicted-vs-measured to be added).

\subsection{Fidelity}
\label{sec:sim-fidelity}

\paraq{On what setting do we demonstrate the simulator's fidelity, and what counts as good?}
We demonstrate fidelity on \textbf{papers100M} at the default \textbf{4$\times$A100}
setting, for both \textbf{SAGE} and \textbf{GAT} at \textbf{$L{=}3$}, comparing the
simulator's predicted per-epoch time against the measured run. The pass criterion is
simple: a prediction \textbf{within 10\%} of measured is considered good.

\begin{table}[h]
\centering
\caption{Simulator fidelity on \textbf{papers100M}, 4$\times$A100, $L{=}3$, hidden 256,
fanout 15. Pass if $|\text{Err}| \le 10\%$. \placeholder{} = to fill.}
\label{tab:sim-fidelity}
\begin{tabular}{l r r r c}
\toprule
Model & Sim (s) & Measured (s) & $|\text{Err}|$ (\%) & within 10\%? \\
\midrule
SAGE & \placeholder & \placeholder & \placeholder & \placeholder \\
GAT  & \placeholder & \placeholder & \placeholder & \placeholder \\
\bottomrule
\end{tabular}
\end{table}

\paraq{What do we expect the fidelity check to show?}
\expected{both SAGE and GAT predictions land within 10\% of measured on papers100M at
$L{=}3$, confirming the compute and data-loading models hold on the large graph at the
default 4-GPU setting.} A miss beyond 10\% would point at a specific cost term---most
likely data-loading, where the on-GPU feature cache must be modelled---rather than at the
approach.

\subsection{Microprofile Cross-Check (nsys)}
\label{sec:microprofile}

\paraq{Does an independent microprofile confirm the cost model's components?}
To verify the cost model at the component level we microprofiled, on one representative
batch, the \emph{fused} full-model forward$+$backward against the simulator's
\emph{per-layer isolated} forward$+$backward, swept over hidden size and model, with an
\texttt{nsys} cross-check, and measured the real pinned-host$\rightarrow$GPU feature gather.

\begin{table}[h]
\centering
\caption{Compute: simulator per-layer vs.\ fused full-model step (products, depth~3,
per-step ms), after the per-layer attribution fix. ratio $=$ sim/fused; last column is
the ratio \emph{before} the fix.}
\label{tab:micro-compute}
\begin{tabular}{l r r r r r}
\toprule
Model & hidden & Fused (ms) & Sim (ms) & ratio & (ratio before) \\
\midrule
SAGE & 64  & 6.02  & 8.66  & 1.44 & 1.46 \\
SAGE & 128 & 5.16  & 7.44  & 1.44 & 2.28 \\
SAGE & 256 & 9.61  & 10.11 & \textbf{1.05} & 2.27 \\
SAGE & 512 & 15.84 & 16.13 & \textbf{1.02} & 3.48 \\
GAT  & 64  & 15.69 & 30.69 & 1.96 & 1.94 \\
GAT  & 128 & 25.05 & 38.86 & 1.55 & 1.58 \\
GAT  & 256 & 40.34 & 52.34 & \textbf{1.30} & 1.56 \\
GAT  & 512 & 79.69 & 92.82 & \textbf{1.16} & 1.72 \\
\bottomrule
\end{tabular}
\end{table}

\paraq{What did the microprofile establish?}
(1)~\textbf{Bandwidth is calibrated to hardware}: NVLink is $300$\,GB/s (\texttt{nvidia-smi
topo}: GPU0--GPU1 are NV12), and the pinned-host$\rightarrow$GPU gather sustains
$\sim$25--26\,GB/s independent of feature size, so \texttt{PCIE\_BW} was corrected
$16\rightarrow25$\,GB/s. (2)~\textbf{Feature loading is linear in feature size}, so the
load/comm terms project by a single scale factor. (3)~\textbf{Compute is GPU-bound GEMM}
(\texttt{nsys}: \texttt{ampere\_sgemm\_*} dominate). (4)~After correcting the
hop$\rightarrow$layer pairing ($L{-}1{-}j$ instead of $-j$), the per-layer estimate
matches the fused step within $\sim$2--5\% for SAGE at hidden $\ge 256$---the large-hidden
regime where the hybrid decision actually matters.

\subsection{Calibration (Step~3a): what it actually required}
\label{sec:sim-calibration}

\paraq{Was calibration just fitting bandwidth constants?}
No---it required \emph{structural} corrections to the cost model as well as fitting
constants, because the uncalibrated simulator under-predicted the measured 4$\times$A100
split time by \textbf{44--66\%} and so always selected $k^\star{=}0$. The ground truth was
obtained \emph{by subtraction} from the measured sweep, with no extra multi-GPU benchmark:
the all-to-all shuffle is not a separate measured column---it is embedded inside the
measured \texttt{forward}/\texttt{backward} (the shuffle autograd runs inside the layer)---so
holding compute fixed, $t^{\text{pcie}}_{\text{fwd+bwd}}-t^{\text{nvlink}}_{\text{fwd+bwd}}$
isolates the real shuffle cost per cell, and the measured NVLink $\text{fwd}{+}\text{bwd}$
gives the per-GPU compute. Both grow with depth/hidden and, on PCIe, the shuffle reaches
$54$--$60\%$ of split time on the deep/large-hidden cells.

\paraq{What were the three corrections?}
(1)~\textbf{A separate all-to-all bandwidth.} The shuffle previously reused the feature-load
\texttt{PCIE\_BW}/\texttt{NVLINK\_BW} (near line rate), but the GPU$\leftrightarrow$GPU
all-to-all pays NCCL and (under PCIe emulation) host-staging overhead, so its \emph{effective}
bandwidth is far lower; fit to the measured shuffle, $\texttt{ALLTOALL\_PCIE\_BW}{\approx}1.95$\,GB/s
(vs.\ the $25$\,GB/s feature-load rate). (2)~\textbf{A missing sampling term.} The model had
\emph{no} neighbour-sampling cost though measured sampling is $1$--$4$\,s; it is now timed on
the real sampler. (3)~\textbf{A per-GPU compute model.} Compute was timed on the \emph{full}
global batch on one GPU, but the real system partitions it across the GPUs; split-layer
compute is now divided by an empirical parallel-efficiency factor
($\texttt{COMPUTE\_PARALLEL\_FACTOR}{\approx}1.7$, fit to the measured per-GPU
$\text{fwd}{+}\text{bwd}$---\emph{not} $\text{num\_devices}{=}4$, because small batches
underutilise the GPU and partition boundaries replicate halo nodes), and DP-layer compute by
$\text{num\_devices}$ (genuinely parallel redundant work).

\begin{table}[h]
\centering
\caption{Calibration fidelity on \textbf{products}, 4$\times$A100, SAGE: the calibrated
simulator's full-split ($k{=}0$) prediction vs.\ the measured split per-epoch time (s).
\emph{Preliminary}: constants fit to these four cells; single measured run; PCIe emulated
(magnitudes directional). The uncalibrated model gave ratios $0.33$--$0.56$.}
\label{tab:sim-calib-products}
\begin{tabular}{l r r r r}
\toprule
Cell (hw, $L$, hid) & Sim $k{=}0$ (s) & Measured split (s) & ratio \\
\midrule
pcie, $L{=}3$, 256  & 4.31  & 4.81  & 0.90 \\
pcie, $L{=}3$, 1024 & 9.84  & 9.41  & 1.05 \\
pcie, $L{=}4$, 256  & 14.27 & 17.74 & 0.80 \\
pcie, $L{=}4$, 1024 & 49.49 & 51.11 & 0.97 \\
\bottomrule
\end{tabular}
\end{table}

\paraq{How honest is this fit?}
The three constants were fit to the four products/SAGE cells above and are not claimed to
transfer unchanged to GAT or other graphs; $\texttt{COMPUTE\_PARALLEL\_FACTOR}$ is a single
compromise value (the true ratio ranges $\sim$1--1.8 across cell sizes). These are therefore
\emph{directional, preliminary} calibrations pending a per-cell refit and the final sweep.

\paraq{Conclusion of this section.}
The simulator predicts a configuration's performance from cheap single-GPU measurements,
and---validated against the DGL baseline and demonstrated within 10\% on papers100M at the
4-GPU setting (plus the component-level microprofile cross-check)---is trustworthy enough
to drive the configuration search of the next section.
